
\documentclass[aps,prd,twocolumn,amsmath,amssymb,longbibliography,nofootinbib]{revtex4-2}

% --- fonts & input (pdfLaTeX preferred by PRD) ---
\usepackage{lmodern}
\usepackage[T1]{fontenc}
\usepackage[utf8]{inputenc}

% --- math & utilities ---
\usepackage{bm}
\usepackage{physics}
\usepackage{siunitx}
\usepackage{graphicx}
\usepackage{xcolor}
\usepackage[final,tracking=true,expansion=false]{microtype}
\usepackage[hidelinks]{hyperref}
\usepackage[nameinlink,capitalize]{cleveref}

\begin{document}

\title{Planck-Scale Consistency of a \texorpdfstring{$10^{-22}\,\mathrm{eV}$}{10^{-22} eV} Bose--Einstein Vacuum Condensate}
\author{Rob Skavberg}
\affiliation{Your Affiliation Here}
\date{January 2026}

\begin{abstract}
We demonstrate that a vacuum composed of ultralight \(10^{-22}\,\mathrm{eV}\) bosons naturally forms a phase-locked Bose--Einstein condensate (BEC) whose macroscopic coherence is directly governed by Planck's constant. This condensate exhibits a bulk modulus \(\beta = \rho_{m} c^{2}\), establishing a mechanical interpretation of the gravitational coupling traditionally appearing in the Einstein--Hilbert action. We show that the quantum pressure, healing length, and coherence scales derived from \(\hbar\) generate macroscopic structure across galaxies, establishing the vacuum as a Planck-consistent quantum medium supporting an emergent gravitational coupling.
\end{abstract}

\maketitle

\section{Introduction}
Planck's constant \(\hbar\) governs quantum fluctuations, coherence, and the phase dynamics of all fields. Traditional formulations treat this quantum structure as distinct from the geometric curvature of general relativity. In the Condensate Link (CL) framework, the vacuum itself is a macroscopic quantum phase, coherent at cosmological scales and capable of producing gravitational behavior as a low-energy mechanical response. Ultralight bosons with mass \(m\sim 10^{-22}\,\mathrm{eV}\) exhibit de Broglie wavelengths of order light-years, leading to significant wavefunction overlap and Bose--Einstein condensation on galactic scales\,\cite{Hui2017ULDM,Marsh2016axion}. These properties allow Planck-scale quantum structure to manifest in macroscopic cosmic behavior.

\section{Planck's Constant and the Emergence of Vacuum Coherence}
Planck's constant sets the scale of the de Broglie wavelength \(\lambda = h/p\), healing length
\begin{equation}
  \xi = \frac{\hbar}{\sqrt{2 m g n}},
\end{equation}
and superfluid stiffness. The Gross--Pitaevskii and Bogoliubov frameworks\,\cite{Gross1961,Pitaevskii1961,Bogoliubov1947,Dalfovo1999,PethickSmith2002} describe the coherent phase, collective excitations, and hydrodynamic limits of BECs. Since \(\hbar\) determines coherence, a vacuum BEC must satisfy Planck-governed phase relations. At the ultralight mass \(10^{-22}\,\mathrm{eV}\), the overlap of wavefunctions is enormous, producing macroscopic coherence.

\section{Why the \texorpdfstring{$10^{-22}$}{1e-22} eV Scale Satisfies Planck Constraints}
\subsection{Coherence on Cosmological Scales}
The Compton wavelength
\begin{equation}
  \lambda_C = \frac{h}{m c}
\end{equation}
for \(m = 10^{-22}\,\mathrm{eV}\) is approximately a light-year, enabling coherent behavior across galactic halos\,\cite{Hui2017ULDM,Marsh2016axion}. Studies of ultralight dark matter (ULDM) show that such particles naturally form condensates whose quantum pressure shapes galaxy structure.

\subsection{Planck-Derived Quantum Pressure}
In the Bogoliubov framework, the quantum pressure scales as
\begin{equation}
  P_{\mathrm q} \propto \frac{\hbar^{2}}{m^{2}} n^{5/3},
\end{equation}
which prevents gravitational collapse in ULDM halos\,\cite{Hui2017ULDM,Marsh2016axion}. This is a direct manifestation of Planck's constant at macroscopic scales.

\section{Mechanical Stiffness and the Einstein--Hilbert Coupling}
The condensate's bulk modulus is
\begin{equation}
  \beta = g\, n^{2},
\end{equation}
while relativistic hydrodynamics (acoustic limit) yields
\begin{equation}
  c^{2} = \frac{\beta}{\rho_{m}}.
\end{equation}
Identifying this \(c\) with the observed phase-velocity limit implies
\begin{equation}
  \beta = \rho_{m} c^{2}.
\end{equation}
The Einstein--Hilbert coupling constant is \(8\pi G / c^{4}\)\,\cite{Einstein1915,EinsteinRelativity1922}. Equating geometric and mechanical stiffness,
\begin{equation}
  \beta = \frac{c^{4}}{8\pi G},
\end{equation}
leads to the emergent identity
\begin{equation}
  G = \frac{c^{2}}{8\pi\, \rho_{m}}.
\end{equation}
Thus, \(G\) arises from condensate density, linking gravitational coupling to Planck-structured quantum matter.

\section{Vacuum as a Phase-Locked Medium}
\subsection{Global Phase Coherence}
The phase \(\theta\) of the condensate is globally coherent because \(\xi\) is extremely large at ultralight mass. This coherence is directly governed by \(\hbar\).

\subsection{Stability}
Quantum pressure balances gravitational stress, producing a stable vacuum medium\,\cite{Hui2017ULDM}. This stabilizing mechanism is a Planck-scale effect manifesting in cosmic structure.

\section{Conclusion}
A vacuum composed of ultralight \(10^{-22}\,\mathrm{eV}\) bosons forms a Bose--Einstein condensate whose coherence, stiffness, and stability arise from Planck's constant. The resulting mechanical identity for \(G\) unifies quantum coherence with gravitational coupling. The CL framework therefore provides a Planck-consistent physical substrate capable of supporting spacetime geometry, light propagation, and gravitational behavior.

\begin{acknowledgments}
We thank colleagues and reviewers for helpful feedback.
\end{acknowledgments}

\bibliographystyle{apsrev4-2}
\bibliography{bib/references}

\end{document}
